\chapter{Contexto do Projeto}
Este capítulo apresenta o contexto real no qual o projeto em desenvolvimento está inserido, caracterizando a área de atuação do parceiro comercial, o perfil do público-alvo, as necessidades operacionais identificadas e a fundamentação conceitual das regras de negócio que determinam o domínio da aplicação web para a confeitaria Rô Bolos. A sua finalidade é estabelecer a compreensão aprofundada dos desafios cotidianos enfrentados pelo estabelecimento antes da especificação e estruturação técnica da solução computacional.

\section{Cenário}

\subsection{O Setor de Confeitaria Artesanal e Food Service no Brasil}
O setor de \textit{food service} e de confeitaria artesanal desempenha um papel de grande relevância na economia nacional, destacando-se como um dos segmentos de maior dinamismo no âmbito das Micro e Pequenas Empresas (MPEs) e do empreendedorismo familiar (SEBRAE, 2023). A produção alimentícia de pequeno porte caracteriza-se por um modelo de produção fortemente baseado em mão de obra artesanal, no qual a personalização dos produtos e a proximidade no atendimento ao cliente constituem os principais diferenciais competitivos frente à produção industrial em larga escala (CHIAVENATO, 2021).

Entretanto, o mercado da confeitaria enfrenta desafios operacionais e econômicos importantes:

\begin{itemize}
    \item A elevada instabilidade nos preços de matérias-primas essenciais (como leite, açúcar, cacau, farinha e derivados) exige uma gestão rigorosa da compra e do consumo de estoque.
    \item A perecibilidade dos insumos e dos produtos acabados impõe prazos de validade curtos, aumentando o risco de perdas financeiras pelos desperdícios.
    \item A falta de regularidade das encomendas demandam planejamento e sintonia entre o recebimento de pedidos, a aquisição de insumos e a escala de produção na cozinha.
\end{itemize}

Em microempreendimentos artesanais, é frequente que a gestão administrativa seja conduzida de maneira informal, com anotações dispersas em cadernos ou trocas de mensagens instantâneas (LAUDON; LAUDON, 2014). Essa ausência de padronização impede o microempreendedor de apurar com exatidão o Custo das Mercadorias Vendidas (CMV), dificultando a precificação dos produtos, o que acaba comprometendo a saúde financeira e a expansão do negócio (PRESSMAN; MAXIM, 2021).

\section{Público-Alvo}

\section{Necessidades Identificadas}

\section{Regras de Negócio}